% !TEX program = pdflatex
\documentclass{article}
\usepackage{iclr2027_conference,times}
\usepackage{hyperref}
\usepackage{url}

\title{Measuring Folded Card Spacers for Model Rails}

\author{Model Rail Desk \\
Paper Model Workshop \\
\texttt{desk@paper-model.example} \\
\And
Card Stock Bench \\
Scale Fixture Room, Paper Model Workshop \\
\texttt{bench@paper-model.example}
}

\begin{document}
\maketitle

\begin{abstract}
A folded card spacer fixes the height of a model rail above its baseboard.
We measured spacers folded from three card stocks in a single test slot and
compared the measured heights with a simple layer model. The model predicts the
height of a folded spacer to within a tenth of a millimetre when the crease
allowance is taken from a two-layer trial. We describe the slot, the measurement
procedure and the remaining sources of error, and we list practical rules for
choosing a card stock for a given rail height.
\end{abstract}

\section{Introduction}\label{sec:intro}
Paper models of railways need many small parts of exact height. Rails,
platform edges and signal bases all sit on spacers, and a spacer that is too
thick by a fraction of a millimetre lifts a rail joint enough to derail a light
wagon. Most builders cut spacers from card and fold them until they reach the
height they want \citep{desk25}. Folding is quick and needs no tools, but the
height of a folded spacer is not simply the sum of the card thicknesses: each
crease adds a small allowance, and the allowance depends on the card and on how
firmly the crease is pressed.

This note reports a bench trial of folded spacers. Our aims were:
\begin{itemize}
\item to measure the height of spacers with two to six layers in one fixed slot;
\item to compare the measurements with a layer model that adds a crease allowance; and
\item to give simple rules for choosing a card stock for a target height.
\end{itemize}
Section~\ref{sec:method} describes the slot and the procedure, Section~\ref{sec:results}
reports the measurements, and Section~\ref{sec:conclusion} summarises the rules.
Earlier notes on layered card \citep{bench24} and on gauge blocks
\citep{scale23} motivated the layer model used here.\footnote{All parts were
cut from the same sheets on the same afternoon, so the moisture content of the
card was the same for every spacer.}

\section{Method}\label{sec:method}
\subsection{The test slot}\label{sec:slot}
The slot was cut into a block of laminated board and lined with a strip of thin
brass so that the card could not dent the walls. A sliding rail of known height
rests on top of the spacer, and the height is read from a scale on the side of
the block. The block was clamped to the bench throughout the trial so that the
scale did not move between readings.

Each spacer was inserted with its folded edge towards the back of the slot. The
rail was lowered onto it and left for ten seconds before the reading was taken,
because a freshly folded spacer springs back slowly. Three readings were taken
for each spacer, removing and reinserting it between readings.

\subsection{Layer model}\label{sec:model}
We model the height $h$ of a spacer with $n$ layers of card of thickness $t$ as
\begin{equation}
h = n\,t + (n-1)\,c ,
\label{eq:height}
\end{equation}
where $c$ is the crease allowance added by each fold. The allowance is found
from a two-layer trial, for which Equation~(\ref{eq:height}) gives
$c = h_2 - 2t$. The same model was used for the gauge block of
\citet{scale23}, where it predicted the height of folded board well for up to
four layers.

\paragraph{Card thickness.}
The thickness $t$ was measured with a dial gauge on ten single sheets of each
stock, and the mean was used in the model. The thin stock measured 0.20\,mm,
the medium stock 0.30\,mm and the thick stock 0.45\,mm, with a spread of less
than 0.01\,mm in every case.

\paragraph{Pressing the crease.}
Every crease was pressed once with the flat of a bone folder along a steel
rule. Pressing twice made the crease allowance smaller but also cracked the
surface of the thick stock, so a single pass was used for all spacers.

\section{Results}\label{sec:results}
Figure~\ref{fig:slot} shows the slot with a four-layer spacer in place, and
Table~\ref{tab:heights} lists the mean measured height of each spacer together
with the height predicted by Equation~(\ref{eq:height}).

\begin{figure}[t]
\centering
\rule{0.8\linewidth}{0.4pt}\\[2pt]
\rule{0.6\linewidth}{18mm}\\[2pt]
\rule{0.8\linewidth}{0.4pt}
\caption{Side view of the test slot with a four-layer spacer under the sliding
rail. The dark block is the spacer; the thin lines mark the slot floor and the
rail.}
\label{fig:slot}
\end{figure}

\begin{table}[t]
\caption{Measured and predicted heights of folded spacers (mm).}
\label{tab:heights}
\centering
\begin{tabular}{lrrr}
\hline
Stock & Layers & Measured & Predicted \\
\hline
Thin & 2 & 0.46 & 0.46 \\
Thin & 4 & 0.95 & 0.98 \\
Medium & 3 & 1.02 & 1.00 \\
Medium & 6 & 2.05 & 2.05 \\
Thick & 4 & 1.98 & 1.92 \\
\hline
\end{tabular}
\end{table}

The model agrees with the measurements to within 0.06\,mm for every spacer. The
largest error is for the thick stock, where the crease allowance grows with the
number of layers: the outer layers must bend around a thicker core, and the
crease opens slightly. For the thin and medium stocks the error does not grow
with the number of layers, which agrees with the note of \citet{glue22} on
thin paper joints.

Repeated readings of the same spacer differed by at most 0.02\,mm. Spacers that
were left in the slot overnight lost about 0.01\,mm of height, which is smaller
than the spread between readings and was ignored.

\section{Conclusion}\label{sec:conclusion}
A folded card spacer can be designed from the card thickness and a crease
allowance measured once for each stock. For heights up to 2\,mm the thin and
medium stocks follow the layer model closely, and they are the better choice
when a rail must sit at an exact height. The thick stock is quicker to fold but
should be checked in the slot when more than three layers are used. A card
ladder \citep{window21} gives a convenient way to check many spacers at once.

\begin{thebibliography}{5}
\bibitem[Card Stock Bench and Model Rail Desk(2024)]{bench24}
Card Stock Bench and Model Rail Desk. 2024.
\newblock Folding tolerances of layered card.
\newblock In \emph{Proceedings of the Bench Meeting on Paper Models}, pages 12--19.

\bibitem[Glue Pot Corner(2022)]{glue22}
Glue Pot Corner. 2022.
\newblock Drying times for thin paper joints.
\newblock \emph{Paper Model Quarterly}, 4(2):33--41.

\bibitem[Model Rail Desk(2025)]{desk25}
Model Rail Desk. 2025.
\newblock \emph{Spacer Card Register}.
\newblock Trial note, Paper Model Workshop.

\bibitem[Scale Fixture Room(2023)]{scale23}
Scale Fixture Room. 2023.
\newblock A gauge block made of folded board.
\newblock Workshop Leaflet~7, Paper Model Workshop.

\bibitem[Window Frame Study(2021)]{window21}
Window Frame Study. 2021.
\newblock Measuring slots with a card ladder.
\newblock Technical memo, Paper Model Workshop.

\end{thebibliography}

\end{document}
